% =============================================================================
%  MAIN ACADEMIC BASE TEMPLATE — BITCOIN HARMONIC TIME MODEL
% =============================================================================
\documentclass[11pt,a4paper,oneside]{article}

% --- Core Packages ---
\usepackage[utf8]{inputenc}
\usepackage[T1]{fontenc}
\usepackage[english]{babel}
\usepackage{amsmath,amsfonts,amssymb,amsthm}
\usepackage{geometry}
\usepackage{graphicx}
\usepackage{booktabs}
\usepackage{tabularx}
\usepackage{siunitx}
\usepackage{float}
\usepackage{hyperref}
\usepackage{microtype}

% --- Page Geometry ---
\geometry{
    a4paper,
    total={160mm,247mm},
    left=25mm,
    top=25mm,
}

% --- SIunitx Configuration ---
\sisetup{
    group-separator={,},
    group-minimum-digits=4
}

% --- Hyperref Setup ---
\hypersetup{
    colorlinks=true,
    linkcolor=blue,
    filecolor=magenta,      
    urlcolor=cyan,
    pdftitle={The Bitcoin Harmonic Time Model},
    pdfpagemode=FullScreen,
}

% --- Mathematical Environments ---
\newtheorem{theorem}{Theorem}
\newtheorem{proposition}{Proposition}
\theoremstyle{definition}
\newtheorem{definition}{Definition}

% =============================================================================
%  METADATA & TITLE CONFIGURATION
% =============================================================================
\title{\textbf{The Bitcoin Harmonic Time Model:\\
Internal Phase Symmetry, Amplitude Power-Laws, and the Institutional ETF Regime Shift}}

\author{\textbf{Romero Quant} \\
\small Independent Quantitative Research \\
\small \texttt{romero.quant@proton.me} \\
\small \href{https://x.com/RomeroQuant}{@RomeroQuant}}

\date{\small June 2026}

% =============================================================================
%  DOCUMENT BODY
% =============================================================================
\begin{document}

\maketitle

\begin{abstract}
This paper formalises the Bitcoin Harmonic Time Model, a two-dimensional quantitative framework that structures Bitcoin market cycles into a discrete dynamic system defined by temporal phase symmetry and compressed amplitude distributions. We demonstrate that the governing clock of the market is not the halving event in isolation, but a stable internal phase architecture where the halving acts as a temporal node, consistently arriving at normalised phase $\bar{\phi}(H) \approx 0.494$ of the expansion phase. Using historical data from three completed cycles, we establish that expansion lasts $1{,}060 \pm 12$ days and contraction $370 \pm 13$ days, yielding a mean period of $1{,}430$ days. Furthermore, we derive an empirical power-law demonstrating that cycle multipliers scale non-linearly with prior drawdown depth ($\hat{\mathcal{M}}(d) = 328 \cdot d^{5.1}$), and introduce a structural adjustment framework to account for the regime shift initiated by the launch of spot US Exchange-Traded Funds (ETFs) in 2024. Finally, we translate this unified model into explicit, dated, and falsifiable calendar hypotheses for the current and subsequent market cycles.
\end{abstract}

\newpage
\tableofcontents
\newpage

% =============================================================================
%  SECTION INTEGRATION (MODULAR INPUTS)
% =============================================================================

% 1. Section 2: Empirical Foundations
% =============================================================================
%  SECTION 2 — EMPIRICAL DATA & TEMPORAL METHODOLOGY
%  Integrate this block immediately after Section 1 in the main document.
% =============================================================================

\section{Empirical Data and Temporal Methodology}

\subsection{Identification of Cycle Extrema}

The empirical foundation of this model rests on the precise identification of
two classes of structural inflection points in Bitcoin's price history: absolute
\textit{cycle troughs} ($T_i$) and absolute \textit{cycle peaks} ($P_i$), where
the subscript $i \in \{1, 2, 3\}$ denotes the cycle index. These inflection
points are defined under a strict methodological protocol designed to minimise
look-ahead bias and subjective interpolation.

\begin{definition}[Cycle Trough]
A \textit{cycle trough} $T_i$ is defined as the date of the global price
minimum recorded on daily closing data between two consecutive cycle peaks
$P_{i-1}$ and $P_i$, where no subsequent retracement to or below that price
level occurs within the same inter-peak interval.
\end{definition}

\begin{definition}[Cycle Peak]
A \textit{cycle peak} $P_i$ is defined as the date of the global all-time high
(ATH) closing price recorded within the corresponding expansion phase, measured
from trough $T_i$ and bounded above by the subsequent trough $T_{i+1}$.
\end{definition}

By construction, each observation window captures one complete
\textit{expansion phase} $[T_i,\, P_i]$ and one complete \textit{contraction
phase} $[P_i,\, T_{i+1}]$, forming the elementary cycle unit
$\mathcal{C}_i = [T_i,\, P_i,\, T_{i+1}]$.

The halving dates $H_i$ are treated as exogenous and protocol-determined
variables. They are taken from the public blockchain record without adjustment:

\begin{itemize}
  \item $H_1$: 9 July 2016 (block 420{,}000; reward reduced from 25 to 12.5 BTC)
  \item $H_2$: 11 May 2020 (block 630{,}000; reward reduced from 12.5 to 6.25 BTC)
  \item $H_3$: 20 April 2024 (block 840{,}000; reward reduced from 6.25 to 3.125 BTC)
\end{itemize}

All duration measures $\Delta_{\phi}$ are expressed in elapsed calendar days,
computed as:

\begin{equation}
  \Delta_{\phi}(A, B) = \lvert \text{date}(B) - \text{date}(A) \rvert
  \quad \text{in days}
\end{equation}

where $A$ and $B$ are any two dated events in the sequence
$\{T_i,\, H_i,\, P_i\}$.

\medskip
\noindent\textit{Note on Cycle 3.} The peak date for the current cycle
($P_3 = 6$ October 2025) is treated as a \textit{working hypothesis} consistent
with the temporal structure derived in this section. Its inclusion in the
averaged parameters is flagged accordingly throughout the analysis, and all
projections derived from it are explicitly labelled as probabilistic. The model
is designed to be falsifiable and self-correcting: if a higher closing price is
established after the submission date of this paper, the cycle extremum and all
downstream projections must be revised.

% -----------------------------------------------------------------------------
\subsection{Cycle Phase Durations}

Table~\ref{tab:phase_durations} reports the duration of the expansion and
contraction phases for each observed cycle. The \textit{expansion duration}
$\Delta^{+}_i$ is measured from trough $T_i$ to peak $P_i$; the
\textit{contraction duration} $\Delta^{-}_i$ is measured from peak $P_i$ to
the subsequent trough $T_{i+1}$.

\begin{table}[H]
  \centering
  \caption{Bitcoin Market Cycle Phase Durations (Calendar Days)}
  \label{tab:phase_durations}
  \renewcommand{\arraystretch}{1.35}
  \begin{tabularx}{\textwidth}{
    l
    S[table-format=4.0]
    S[table-format=3.0]
  }
    \toprule
    \textbf{Cycle} &
    \multicolumn{1}{c}{\textbf{Expansion Phase} $\Delta^{+}_i$ (days)} &
    \multicolumn{1}{c}{\textbf{Contraction Phase} $\Delta^{-}_i$ (days)} \\
    \midrule
    \textbf{Cycle 1}
      & {1068}
      & {363} \\
    \quad Trough $T_1$: 14 Jan 2015
      & {}
      & {} \\
    \quad Peak $P_1$: 17 Dec 2017
      & {}
      & {} \\
    \quad Trough $T_2$: 15 Dec 2018
      & {}
      & {} \\[4pt]
    \textbf{Cycle 2}
      & {1061}
      & {376} \\
    \quad Trough $T_2$: 15 Dec 2018
      & {}
      & {} \\
    \quad Peak $P_2$: 10 Nov 2021
      & {}
      & {} \\
    \quad Trough $T_3$: 21 Nov 2022
      & {}
      & {} \\[4pt]
    \textbf{Cycle 3} \textsuperscript{\dag}
      & {1050}
      & {---} \\
    \quad Trough $T_3$: 21 Nov 2022
      & {}
      & {} \\
    \quad Peak $P_3$: 6 Oct 2025 \textsuperscript{\dag}
      & {}
      & {} \\
    \midrule
    \textbf{Average} $\bar{\Delta}$
      & \textbf{1060}
      & \textbf{370} \\
    \quad Range $[\min, \max]$
      & {[1050,\; 1068]}
      & {[363,\; 376]} \\
    \quad Inter-cycle deviation $\sigma_{\max}$
      & {18 days}
      & {13 days} \\
    \bottomrule
  \end{tabularx}
  \medskip
  \begin{minipage}{\textwidth}
    \small
    \textsuperscript{\dag} Cycle 3 peak ($P_3$) is a working hypothesis. The
    contraction phase $\Delta^{-}_3$ is therefore excluded from the current
    average computation and will be incorporated upon confirmation of the cycle
    extremum.
  \end{minipage}
\end{table}

The arithmetic mean of the expansion phase across three cycles yields
$\bar{\Delta}^{+} \approx 1{,}060$ days, with a maximum inter-cycle deviation
of $\sigma_{\max}^{+} = 18$ days---representing a relative dispersion of
approximately $1.7\%$ around the mean. The contraction phase, measured over
two complete observations, yields $\bar{\Delta}^{-} \approx 370$ days with a
maximum deviation of $\sigma_{\max}^{-} = 13$ days, a relative dispersion of
$3.5\%$.

% -----------------------------------------------------------------------------
\subsection{Internal Phase Symmetry: The Halving as Temporal Anchor}

The most structurally significant finding of this empirical analysis emerges
when the expansion phase is decomposed at the halving event $H_i$. Rather than
treating the halving as the origin of the cycle---as conventional
halving-centric models implicitly do---we position it as an internal marker
that divides the expansion phase $[T_i, P_i]$ into two measurable sub-intervals:

\begin{align}
  \alpha_i &= \Delta_{\phi}(T_i,\; H_i) \quad \text{(pre-halving segment)} \\
  \beta_i  &= \Delta_{\phi}(H_i,\; P_i) \quad \text{(post-halving segment)}
\end{align}

such that, by construction, $\alpha_i + \beta_i = \Delta^{+}_i$.

\begin{hypothesis}[Internal Phase Symmetry]
  The halving event $H_i$ functions as a \textit{temporal anchor} within the
  expansion phase: the sub-intervals $\alpha_i$ and $\beta_i$ are quasi-equal
  and exhibit inter-cycle stability that is statistically independent of
  price-level dynamics.
\end{hypothesis}

Table~\ref{tab:phase_symmetry} reports the measured values of $\alpha_i$ and
$\beta_i$ for each observed cycle.

\begin{table}[H]
  \centering
  \caption{Internal Phase Symmetry --- Expansion Phase Decomposed at the Halving}
  \label{tab:phase_symmetry}
  \renewcommand{\arraystretch}{1.35}
  \begin{tabularx}{\textwidth}{
    l
    c
    S[table-format=3.0]
    S[table-format=3.0]
    S[table-format=4.0]
  }
    \toprule
    \textbf{Cycle} &
    \textbf{Halving Date} ($H_i$) &
    \multicolumn{1}{c}{$\alpha_i$ (days)} &
    \multicolumn{1}{c}{$\beta_i$ (days)} &
    \multicolumn{1}{c}{$\alpha_i + \beta_i$ (days)} \\
    \midrule
    Cycle 1 (2015--2017) & 9 Jul 2016  & 542 & 526 & 1068 \\
    Cycle 2 (2018--2021) & 11 May 2020 & 513 & 548 & 1061 \\
    Cycle 3 (2022--2025)\textsuperscript{\dag}
                          & 20 Apr 2024 & 515 & 534 & 1050 \\
    \midrule
    \textbf{Average} $\bar{\alpha},\; \bar{\beta}$
      & ---
      & \textbf{523}
      & \textbf{536}
      & \textbf{1059} \\
    \quad Max.\ deviation $\sigma_{\max}$
      & ---
      & {19 days}
      & {12 days}
      & {---} \\
    \bottomrule
  \end{tabularx}
  \medskip
  \begin{minipage}{\textwidth}
    \small
    \textsuperscript{\dag} Cycle 3 values are derived under the hypothesis
    $P_3 = 6$ October 2025.
  \end{minipage}
\end{table}

The data reveal a bilateral stability that is difficult to attribute to
stochastic coincidence. The pre-halving segment $\alpha_i$ ranges from 513 to
542 days, yielding a mean of $\bar{\alpha} = 523$ days and a maximum deviation
of $\sigma_{\max}^{\alpha} = 19$ days ($\approx 3.6\%$ of the mean). The
post-halving segment $\beta_i$ ranges from 526 to 548 days, yielding a mean of
$\bar{\beta} = 536$ days with a maximum deviation of $\sigma_{\max}^{\beta}=
12$ days ($\approx 2.2\%$ of the mean). Taken together, the sum $\bar{\alpha}
+ \bar{\beta} = 1{,}059$ days is consistent to within one day of the
independently computed expansion mean $\bar{\Delta}^{+} = 1{,}060$ days,
providing internal cross-validation of the measurement methodology.

The asymmetry between $\bar{\alpha}$ and $\bar{\beta}$ is also noteworthy: at
13 days, it indicates that the halving does not bisect the expansion phase
exactly at its midpoint, but rather occurs \textit{slightly before} it. In
other words, the system reaches its peak approximately 13 days \textit{later}
in the post-halving segment than what a perfectly symmetric division would
predict. This mild asymmetry is stable across all three cycles and may reflect
the lagged transmission mechanism through which supply reduction propagates
into speculative price discovery.

\begin{figure}[H]
    \centering
    \includegraphics[width=1.0\textwidth]{clock.png}
    \caption{\textbf{The Bitcoin Harmonic Clock:} Visualisation of the internal phase architecture across Cycles 1, 2, and 3. The diagram demonstrates the structural invariance of the expansion phase, positioning the halving event not as an origin, but as an internal temporal node fixed at $\overline{\phi}(H) \approx 0.494$.}
    \label{fig:harmonic_clock}
\end{figure}

% -----------------------------------------------------------------------------
\subsection{Statistical Characterisation of Temporal Stability}

The preceding measurements permit a formal statistical characterisation of the
temporal regularity of Bitcoin's market structure. Let us define the
\textit{temporal stability index} $\mathcal{S}(\phi)$ of a phase $\phi$ as:

\begin{equation}
  \mathcal{S}(\phi) = 1 - \frac{\sigma_{\max}(\phi)}{\bar{\Delta}(\phi)}
\end{equation}

where $\sigma_{\max}(\phi)$ is the maximum observed deviation from the phase
mean and $\bar{\Delta}(\phi)$ is the phase mean duration. A value of
$\mathcal{S}(\phi) = 1$ would denote perfect invariance; values approaching
zero indicate high temporal dispersion.

Applying this definition to the empirically observed phases:

\begin{align}
  \mathcal{S}(\Delta^{+}) &= 1 - \frac{18}{1060} \approx 0.983 \\[4pt]
  \mathcal{S}(\Delta^{-}) &= 1 - \frac{13}{370}  \approx 0.965 \\[4pt]
  \mathcal{S}(\alpha)      &= 1 - \frac{19}{523}  \approx 0.964 \\[4pt]
  \mathcal{S}(\beta)       &= 1 - \frac{12}{536}  \approx 0.978
\end{align}

All four phase parameters exhibit temporal stability indices in excess of
$0.96$. This is a remarkable result for a market that is widely characterised
as chaotic and subject to reflexive price dynamics. Most notably, the full
expansion phase $\Delta^{+}$ returns the highest stability index of the four,
$\mathcal{S}(\Delta^{+}) \approx 0.983$, implying that the bullish trajectory
from cycle trough to cycle peak is the most temporally invariant structural
unit in the system.

\begin{proposition}[Differential Temporal Stability]
  The expansion phase of Bitcoin market cycles is more temporally stable than
  the contraction phase: $\mathcal{S}(\Delta^{+}) > \mathcal{S}(\Delta^{-})$.
  Equivalently, the maximum inter-cycle deviation of the expansion phase
  (18 days over a mean of 1{,}060 days) is proportionally smaller than that of
  the contraction phase (13 days over a mean of 370 days).
\end{proposition}

This proposition carries a significant analytical implication. The prevailing
market intuition often treats bear markets as more structurally predictable
than bull markets---a view shaped by the observation that fear-driven
capitulation tends to be temporally compressed and abrupt. The data presented
here challenge that intuition at the cycle level: while individual contraction
events may feel decisive, their total duration is less invariant across cycles
than that of the expansion phases that precede them. Put differently,
\textit{the system's internal clock runs more precisely during expansions than
during contractions.}

This asymmetry may reflect a fundamental difference in the dynamics governing
the two phases. The expansion phase is characterised by a progressive and
relatively orderly accumulation of demand---a process that, by its nature,
unfolds over extended timeframes and is anchored by the halving's structural
signal. The contraction phase, by contrast, is governed by the interaction
between forced selling, liquidity withdrawal, and sentiment deterioration:
forces that are inherently more sensitive to exogenous shocks and therefore
more variable in their temporal manifestation.

Critically, the observation that the expansion phase post-ETF approval (Cycle 3)
continues to conform to the historical temporal envelope---with a duration of
$\Delta^{+}_3 = 1{,}050$ days, only 10 days below the prior minimum of
1{,}061 days---provides preliminary evidence that the temporal structure of
the model is robust to the structural regime change introduced by institutional
participation. The amplitude parameters, as we discuss in the following
sections, do not appear to enjoy the same degree of cross-regime stability.

% =============================================================================
%  END OF SECTION 2
% =============================================================================


% 2. Sections 3 & 4: Phase Formalisation and Law of Amplitude
% =============================================================================
%  SECTION 3 — MATHEMATICAL FORMALISATION OF INTERNAL PHASE SYMMETRY
%  SECTION 4 — THE LAW OF AMPLITUDE AND STRUCTURAL REGIME ADJUSTMENTS
%  Integrate immediately after Section 2 in the main document.
% =============================================================================


% =============================================================================
\section{Mathematical Formalisation of Internal Phase Symmetry}
% =============================================================================

\subsection{Bitcoin Cycles as a Damped Harmonic System}

The empirical regularities documented in Section~2 motivate a formal analogical
framework borrowed from the theory of dynamic systems. We propose that the
observable structure of Bitcoin market cycles admits a tractable representation
as a \textit{discrete damped harmonic oscillator}: a system characterised by
near-constant periodicity, a stable internal phase relationship, and
monotonically decreasing amplitude across successive cycles.

In continuous-time mechanics, a damped harmonic oscillator obeys the
second-order linear differential equation:

\begin{equation}
  \ddot{x}(t) + 2\zeta\omega_0\,\dot{x}(t) + \omega_0^2\,x(t) = 0
  \label{eq:dho}
\end{equation}

where $x(t)$ denotes the displacement from equilibrium, $\omega_0$ the
natural angular frequency of the undamped system, and $\zeta \in (0,1)$ the
damping ratio. For $\zeta \in (0,1)$---the underdamped regime---the general
solution takes the form:

\begin{equation}
  x(t) = A_0\, e^{-\zeta\omega_0 t}
           \cos\!\bigl(\omega_d t + \varphi_0\bigr)
  \label{eq:dho_solution}
\end{equation}

where $\omega_d = \omega_0\sqrt{1 - \zeta^2}$ is the damped angular frequency,
$A_0$ is the initial amplitude, and $\varphi_0$ is the initial phase offset.
The solution describes oscillations of progressively decreasing amplitude
enveloped by the exponential decay factor $e^{-\zeta\omega_0 t}$.

We do not claim that equation~\eqref{eq:dho} governs Bitcoin's price path in
any literal mechanical sense. Markets are reflexive, non-linear systems, and
price dynamics are not solutions to differential equations with fixed
coefficients. The value of the analogy lies in its structural
correspondence: (i) a measurable and nearly constant period $T$; (ii) a
well-defined phase relationship pinned by an internal reference event; and
(iii) a monotonically decreasing amplitude envelope. All three properties are
empirically supported by the data in Section~2 for the temporal dimension and
are developed for the amplitude dimension in Section~4.

\subsection{Cycle Period and Angular Frequency}

The total cycle period $T_{\text{cycle}}$ is defined as the elapsed time
between two consecutive troughs:

\begin{equation}
  T_{\text{cycle}} = \Delta^{+}_i + \Delta^{-}_i
  \label{eq:period}
\end{equation}

Using the empirically derived averages $\bar{\Delta}^{+} \approx 1{,}060$ days
and $\bar{\Delta}^{-} \approx 370$ days:

\begin{equation}
  \bar{T}_{\text{cycle}} \approx 1{,}060 + 370 = 1{,}430 \text{ days}
  \label{eq:period_value}
\end{equation}

The corresponding mean angular frequency $\bar{\omega}$, expressed in
radians per day, is:

\begin{equation}
  \bar{\omega} = \frac{2\pi}{\bar{T}_{\text{cycle}}}
               = \frac{2\pi}{1{,}430}
               \approx 4.394 \times 10^{-3}\ \text{rad}\cdot\text{day}^{-1}
  \label{eq:omega}
\end{equation}

\subsection{Phase Definition and the Halving as a Temporal Node}

Within each cycle $\mathcal{C}_i = [T_i,\, P_i,\, T_{i+1}]$, we define the
instantaneous cycle phase $\phi(t)$ as the normalised position of any date $t$
within the expansion interval $[T_i, P_i]$:

\begin{equation}
  \phi(t) = \frac{t - T_i}{\Delta^{+}_i}, \quad \phi \in [0,\, 1]
  \label{eq:phase_norm}
\end{equation}

Under this normalisation, $\phi(T_i) = 0$ at the cycle trough and
$\phi(P_i) = 1$ at the cycle peak. The phase at the halving event $H_i$ is
therefore:

\begin{equation}
  \phi(H_i) = \frac{\alpha_i}{\Delta^{+}_i}
             = \frac{\alpha_i}{\alpha_i + \beta_i}
  \label{eq:phase_halving}
\end{equation}

Substituting the empirical averages $\bar{\alpha} = 523$ days and
$\bar{\beta} = 536$ days:

\begin{equation}
  \bar{\phi}(H) = \frac{523}{523 + 536}
                = \frac{523}{1{,}059}
                \approx 0.4939
  \label{eq:phase_halving_value}
\end{equation}

The halving consistently occurs at approximately $\bar{\phi}(H) \approx 0.494$
of the expansion phase---marginally before the mathematical midpoint
$\phi = 0.500$. We define the \textit{phase displacement} $\delta_\phi$ as the
deviation of the halving from perfect bisection:

\begin{equation}
  \delta_\phi = \bar{\phi}(H) - 0.500 = 0.494 - 0.500 = -0.006
  \label{eq:phase_displacement}
\end{equation}

This displacement of $-0.006$ in normalised phase corresponds to approximately
$-0.006 \times 1{,}059 \approx -6$ calendar days: the halving arrives,
on average, six days earlier than the geometric midpoint of the expansion. The
sign is negative because the system's peak lies slightly further in time from
the halving than the trough does---a structural asymmetry we interpret as the
lagged propagation of the supply shock through market sentiment and price
discovery mechanisms.

The observed inter-cycle values of $\phi(H_i)$ are:

\begin{align}
  \phi_1(H_1) &= \frac{542}{1068} \approx 0.508 \\[4pt]
  \phi_2(H_2) &= \frac{513}{1061} \approx 0.484 \\[4pt]
  \phi_3(H_3) &= \frac{515}{1050} \approx 0.490
\end{align}

The range of $\phi(H_i)$ is $[0.484,\; 0.508]$, a span of 0.024 in normalised
phase, or approximately 25 calendar days when projected onto the mean expansion
duration. This near-invariance of the halving's phase position across three
independent cycles---each occurring in markedly different macroeconomic and
structural environments---constitutes the strongest quantitative argument for
treating the halving as a \textit{temporal node}: a fixed point in the phase
space of the system, rather than an event whose timing relative to the cycle
is coincidental.

\begin{definition}[Temporal Node]
  The halving event $H_i$ is a \textit{temporal node} of cycle $\mathcal{C}_i$
  if and only if its normalised phase $\phi(H_i)$ is invariant across cycles
  up to a maximum deviation of $\varepsilon_\phi$, where $\varepsilon_\phi$ is
  small relative to unity. Empirically, $\varepsilon_\phi \approx 0.012$ (one
  half of the observed range), satisfying the condition $\varepsilon_\phi \ll 1$.
\end{definition}

\subsection{Implications for Forward Projection}

The phase-anchoring property of the halving has a direct and practically
consequential implication: knowing the date of the next halving $H_{i+1}$
and an estimate of the trough $T_{i+1}$, it is possible to construct a
probabilistic interval for the next peak $P_{i+1}$ without reference to any
price-level information. Specifically:

\begin{equation}
  \hat{P}_{i+1} = H_{i+1} + \bar{\beta} \pm \sigma_{\max}^{\beta}
                = H_{i+1} + 536 \pm 12 \ \text{ days}
  \label{eq:peak_projection}
\end{equation}

This formulation renders the temporal projection \textit{falsifiable} in a
strictly Popperian sense: it generates a dated prediction interval that can
be confirmed or refuted by subsequent market data, independently of any
assumptions about price magnitude. The projection for the current cycle is
developed in full in Section~6.


% =============================================================================
\section{The Law of Amplitude and Structural Regime Adjustments}
% =============================================================================

\subsection{Empirical Amplitude Data}

The second measurable degree of freedom in the harmonic model is the
\textit{amplitude} of each cycle: the magnitude of the price return from
trough to peak. We characterise this dimension through two complementary
empirical series. The first is the \textit{cycle multiplier} $\mathcal{M}_i$,
defined as the ratio of the peak price to the trough price within cycle $i$:

\begin{equation}
  \mathcal{M}_i = \frac{P^{\$}_i}{T^{\$}_i}
  \label{eq:multiplier}
\end{equation}

The second is the \textit{prior drawdown} $d_i$, defined as the proportional
decline from the preceding cycle peak $P^{\$}_{i-1}$ to the current trough
$T^{\$}_i$:

\begin{equation}
  d_i = \frac{P^{\$}_{i-1} - T^{\$}_i}{P^{\$}_{i-1}}
      = 1 - \frac{T^{\$}_i}{P^{\$}_{i-1}}, \quad d_i \in (0,\, 1)
  \label{eq:drawdown}
\end{equation}

Table~\ref{tab:amplitude} reports the empirical values of $\mathcal{M}_i$ and
$d_i$ for the three observed cycles.

\begin{table}[H]
  \centering
  \caption{Bitcoin Cycle Amplitude Parameters: Multipliers and Prior Drawdowns}
  \label{tab:amplitude}
  \renewcommand{\arraystretch}{1.35}
  \begin{tabularx}{\textwidth}{l S[table-format=2.2] S[table-format=1.2]
                                 S[table-format=2.1] c}
    \toprule
    \textbf{Cycle} &
    \multicolumn{1}{c}{$\mathcal{M}_i$ (Trough$\to$Peak)} &
    \multicolumn{1}{c}{$d_i$ (Prior drawdown)} &
    \multicolumn{1}{c}{Return (\%)} &
    \textbf{Regime} \\
    \midrule
    Cycle 1 (2015--2017) & 130.40 & 0.83 & {+12{,}940\%} & Retail \\
    Cycle 2 (2018--2021) &  20.44 & 0.59 & {+1{,}944\%}  & Retail \\
    Cycle 3 (2022--2025)\textsuperscript{\dag}
                          &   8.07 & 0.48 & {+707\%}      & Transitional \\
    \midrule
    \textit{Projected Cycle 4} & \textit{2.5--4.0} & \textit{0.30--0.40?}
      & \textit{+150--300\%} & \textit{Institutional} \\
    \bottomrule
  \end{tabularx}
  \medskip
  \begin{minipage}{\textwidth}
    \small
    \textsuperscript{\dag} Cycle 3 amplitude computed under the hypothesis
    $P_3 = 6$ October 2025. Cycle 4 projections are model-derived estimates
    subject to the regime-adjustment discussed in \S\ref{sec:etf_break}.
  \end{minipage}
\end{table}

Two secular trends are immediately apparent from Table~\ref{tab:amplitude}.
First, the cycle multiplier $\mathcal{M}_i$ declines sharply and
monotonically across successive cycles: from $130.4\times$ to $20.4\times$
to $8.1\times$. Second, the prior drawdown $d_i$ likewise declines: from $83\%$
to $59\%$ to $48\%$. The co-movement of these two series suggests that the
amplitude of each recovery is not independent of the severity of the preceding
contraction---a relationship that we formalise in the following subsection.

\subsection{The Power-Law of Amplitude: Empirical Derivation}

Inspection of the data in Table~\ref{tab:amplitude} suggests that the cycle
multiplier $\mathcal{M}_i$ is functionally related to the prior drawdown $d_i$
through a power-law specification of the form:

\begin{equation}
  \mathcal{M}(d) = \kappa \cdot d^{\,\gamma}
  \label{eq:amplitude_law}
\end{equation}

where $\kappa > 0$ is a scaling coefficient and $\gamma > 0$ is the elasticity
exponent. Fitting equation~\eqref{eq:amplitude_law} to the three observed
data points $\{(d_i, \mathcal{M}_i)\}_{i=1}^{3}$ via least-squares regression
in log-linear space---that is, regressing $\ln\mathcal{M}$ on $\ln d$---yields
the following empirical estimates:

\begin{equation}
  \boxed{\ \hat{\mathcal{M}}(d) \approx 328 \cdot d^{\,5.1}\ }
  \label{eq:amplitude_fit}
\end{equation}

with $\hat{\kappa} = 328$ and $\hat{\gamma} = 5.1$.

The high value of the exponent $\hat{\gamma} = 5.1$ reflects the extreme
non-linearity of the relationship: small reductions in drawdown depth produce
disproportionately large compressions in the subsequent multiplier. This is
the mathematical expression of the \textit{Law of Diminishing Returns} that
characterises maturing asset classes. Table~\ref{tab:amplitude_fit} reports
the in-sample fit of the model against observed values.

\begin{table}[H]
  \centering
  \caption{In-Sample Fit of the Power-Law Amplitude Model
           $\hat{\mathcal{M}}(d) = 328 \cdot d^{5.1}$}
  \label{tab:amplitude_fit}
  \renewcommand{\arraystretch}{1.35}
  \begin{tabularx}{\textwidth}{l
    S[table-format=1.2]
    S[table-format=3.1]
    S[table-format=3.1]
    S[table-format=2.1]}
    \toprule
    \textbf{Cycle} &
    \multicolumn{1}{c}{$d_i$} &
    \multicolumn{1}{c}{$\hat{\mathcal{M}}(d_i)$ (Model)} &
    \multicolumn{1}{c}{$\mathcal{M}_i$ (Observed)} &
    \multicolumn{1}{c}{Error (\%)} \\
    \midrule
    Cycle 1 & 0.83 & 126.0 & 130.4 & {-3.4\%} \\
    Cycle 2 & 0.59 &  22.0 &  20.4 & {+7.8\%} \\
    Cycle 3 & 0.48 &   7.7 &   8.1 & {-4.9\%} \\
    \bottomrule
  \end{tabularx}
\end{table}

The model reproduces the three observed multipliers with errors of less than
$8\%$, a precision that is considered high for a two-parameter empirical model
fitted to three data points across markets separated by a decade and orders of
magnitude in capitalisation. The direction of the errors is unsystematic
(negative, positive, negative), suggesting the absence of a structural bias in
the fit and consistent with random measurement or timing uncertainty rather
than model misspecification.

\begin{figure}[H]
    \centering
    \includegraphics[width=0.9\textwidth]{law_of_amplitude.png}
    \caption{\textbf{The Law of Amplitude and Regime Shift:} Empirical power-law trajectory mapped as a function of prior drawdown severity ($\hat{\mathcal{M}}(d) = 328 \cdot d^{5.1}$). The yellow boundary highlights the projected structural break and volatility dampening induced by the post-2024 spot institutional ETF regime change.}
    \label{fig:law_of_amplitude}
\end{figure}

\subsection{The Law of Diminishing Returns: Structural Interpretation}

The power-law relationship in equation~\eqref{eq:amplitude_fit} is not
merely a curve-fitting exercise. It admits a structural interpretation grounded
in the economics of asset maturation. We identify three reinforcing mechanisms
that produce the observed amplitude compression:

\textbf{Mechanism I: Market Capitalisation Scaling.} As Bitcoin's market
capitalisation increases by an order of magnitude with each cycle, the absolute
capital inflow required to achieve a given price multiplier grows
proportionally. A $100\times$ return from a \$1B market cap requires \$100B
in net new capital; the same return from a \$100B capitalisation requires
\$10T---a figure that exceeds the investable universe of most asset classes. The
\textit{reflexive compression} of multipliers as capitalisation grows is
therefore a near-mathematical inevitability rather than a contingent market
phenomenon.

\textbf{Mechanism II: Drawdown Compression as a Maturity Indicator.} The
monotonic decline in $d_i$ from $0.83$ to $0.59$ to $0.48$ is itself a
manifestation of structural maturation. As the asset develops deeper spot and
derivatives liquidity, a larger and more diversified investor base, and
greater macro-level awareness, the conditions for absolute capitulation---which
characterised the $-83\%$ drawdown of 2014--2015---become progressively less
tenable. Tighter drawdowns imply, via equation~\eqref{eq:amplitude_law} and
the large exponent $\hat{\gamma}=5.1$, dramatically lower subsequent
multipliers. In this sense, the drawdown depth $d_i$ functions as an implicit
\textit{maturity proxy} for the asset at the time of cycle trough.

\textbf{Mechanism III: Reflexive Expectation Dampening.} With each completed
cycle, the universe of potential Bitcoin investors who have not yet participated
shrinks. As the narrative of extraordinary returns becomes part of collective
financial memory, forward-looking participants price in a portion of the
expected cycle gain during the accumulation phase itself, reducing the
magnitude of the subsequent re-rating. This mechanism---analogous to the
well-documented reduction in the post-earnings announcement drift as analyst
coverage increases---produces a secular compression in realised multipliers
independently of capitalisation effects.

\subsection{The Structural Break: Spot ETF Approval as a Regime Shift}
\label{sec:etf_break}

The preceding mechanisms describe endogenous sources of amplitude compression
that would have been operative regardless of the regulatory and product
landscape. Beginning in January 2024, however, a qualitatively distinct and
exogenous force entered the system: the approval and large-scale adoption of
spot Bitcoin Exchange-Traded Funds in the United States and subsequently in
other major jurisdictions.

The ETF channel introduces capital flows with properties that are
fundamentally different from those of prior demand sources:

\begin{enumerate}
  \item \textbf{Passivity and mandate-driven allocation.} A substantial portion
  of ETF inflows originates from multi-asset portfolio rebalancing mandates
  rather than from discretionary conviction. This capital is, by construction,
  less sensitive to price levels and sentiment cycles than retail or
  speculative institutional capital.

  \item \textbf{Continuous and symmetric flow structure.} Unlike retail
  participants, who tend to concentrate purchases near cycle peaks and sales
  near troughs, mandate-driven ETF flows exhibit a more uniform distribution
  across time---dampening both the velocity of price appreciation during
  expansion phases and the velocity of decline during contraction phases.

  \item \textbf{Reduction of uninvested addressable market.} The ETF structure
  makes Bitcoin accessible to a class of capital---pension funds, endowments,
  wealth management platforms---that was previously excluded by custody,
  compliance, or fiduciary constraints. As this capital enters progressively
  during and after Cycle 3, the pool of ``yet-to-participate'' institutional
  demand diminishes, further constraining future multiplier potential.
\end{enumerate}

The net structural effect of ETF participation can be decomposed into two
directional impacts on the amplitude parameters:

\begin{equation}
  \Delta d_i^{\,\text{ETF}} < 0 \quad \Longrightarrow \quad
  \Delta \mathcal{M}_i^{\,\text{ETF}} < 0
  \label{eq:etf_effect}
\end{equation}

That is: ETF participation \textit{compresses the drawdown} $d_i$ below what
endogenous mechanisms would produce, which in turn---through the power-law
relationship of equation~\eqref{eq:amplitude_fit}---\textit{further compresses
the subsequent multiplier} $\mathcal{M}_i$. The two effects are not additive
but multiplicatively linked via the functional form of the amplitude law.

\subsection{Forward Amplitude Projection for Cycle 4}

Under the combined action of endogenous maturation and the ETF-driven
structural break, we project the prior drawdown for Cycle 4 in the range
$d_4 \in [0.30,\; 0.40]$. This range reflects:

\begin{itemize}
  \item The secular trend of declining drawdowns: $0.83 \to 0.59 \to 0.48$,
  suggesting a continuation toward the $[0.30, 0.40]$ interval.
  \item The additional stabilising force of sustained ETF demand during
  contraction phases, which introduces a structural price floor absent in
  prior cycles.
  \item Uncertainty about the magnitude of ETF-driven support, which precludes
  a point estimate and warrants the use of an interval.
\end{itemize}

Applying the power-law model of equation~\eqref{eq:amplitude_fit} to this
drawdown range, whilst acknowledging that the model was calibrated on
pre-ETF data and may therefore \textit{overestimate} the multiplier under the
new regime:

\begin{align}
  \hat{\mathcal{M}}(0.40) &= 328 \cdot (0.40)^{5.1}
    \approx 328 \cdot 0.00926 \approx 3.04 \label{eq:m_upper} \\[4pt]
  \hat{\mathcal{M}}(0.30) &= 328 \cdot (0.30)^{5.1}
    \approx 328 \cdot 0.00218 \approx 0.71 \label{eq:m_lower}
\end{align}

The lower bound produced by equation~\eqref{eq:m_lower} ($d_4 = 0.30$,
$\hat{\mathcal{M}} \approx 0.71$) implies a peak price below the trough price,
which is structurally inconsistent with a continuation of the harmonic cycle.
We interpret this as evidence that the purely mechanical extrapolation of the
power-law model breaks down at low drawdown levels---a regime where the
ETF-driven price floor prevents the kind of absolute capitulation that the
model was fitted to capture.

We therefore propose a \textit{regime-adjusted amplitude interval} for Cycle 4
that combines the mechanical model output for $d_4 \in [0.35, 0.40]$ with a
qualitative adjustment for ETF structural support:

\begin{equation}
  \mathcal{M}_4^{\text{adj}} \in [2.5,\; 4.0]
  \label{eq:m4_adjusted}
\end{equation}

This interval implies a percentage return from trough to peak of approximately
$+150\%$ to $+300\%$---a range that reflects the transition from a
hypergrowth asset to a maturing store-of-value instrument with deep
institutional participation. Critically, this projection is derived from the
structural and mechanical properties of the model rather than from price
extrapolation, and it should be treated as a probabilistic envelope rather
than a point forecast.

\begin{proposition}[Amplitude Convergence]
  Under continued asset maturation and sustained institutional ETF
  participation, the cycle multiplier $\mathcal{M}_i$ is expected to converge
  asymptotically toward a lower bound consistent with the return profile of
  other macro assets. The power-law relationship
  $\hat{\mathcal{M}}(d) = 328 \cdot d^{5.1}$ captures this convergence
  endogenously through the compression of prior drawdowns, but requires a
  regime adjustment for cycles in which ETF-driven demand materially
  alters the drawdown distribution relative to the pre-ETF calibration sample.
\end{proposition}

% =============================================================================
%  END OF SECTIONS 3 & 4
% =============================================================================


% 3. Sections 5, 6 & 7: Synthesis, Projections and Limitations
% =============================================================================
%  SECTION 5 — SYNTHESIS OF THE UNIFIED HARMONIC MODEL
%  SECTION 6 — FORWARD PROJECTIONS AS FALSIFIABLE HYPOTHESES
%  SECTION 7 — CONCLUDING REMARKS AND LIMITATIONS
%
%  Integrate immediately after Sections 3 & 4.
%  This block does NOT contain \end{document} — that remains in the base template.
% =============================================================================


% =============================================================================
\section{Synthesis of the Unified Harmonic Model}
% =============================================================================

\subsection{Two Independent Dimensions of a Single Dynamic System}

The preceding sections have developed two analytically distinct but
structurally interdependent dimensions of Bitcoin's market cycle. Section~2
and Section~3 established the \textit{temporal dimension}: the expansion and
contraction phases exhibit near-invariant durations of $\bar{\Delta}^{+}
\approx 1{,}060$ days and $\bar{\Delta}^{-} \approx 370$ days, with the
halving functioning as a stable temporal node at phase $\bar{\phi}(H) \approx
0.494$ of the expansion. Section~4 established the \textit{amplitude
dimension}: cycle multipliers follow an empirical power-law
$\hat{\mathcal{M}}(d) = 328 \cdot d^{5.1}$ as a function of prior drawdown
depth, subject to a structural downward adjustment under the ETF regime.

The central claim of this synthesis is that these two dimensions are
\textit{quasi-orthogonal} in their dynamics yet \textit{jointly necessary} for
a complete characterisation of the system. Formally, we represent the state of
the Bitcoin market at any point in time $t$ as a two-dimensional vector
$\mathbf{S}(t)$:

\begin{equation}
  \mathbf{S}(t) = \bigl(\phi(t),\; \mathcal{A}(t)\bigr)
  \label{eq:state_vector}
\end{equation}

where $\phi(t) \in [0, 1]$ is the normalised cycle phase defined in
equation~\eqref{eq:phase_norm} and $\mathcal{A}(t) \geq 0$ is the
instantaneous amplitude level, normalised to the trough price
$T^{\$}_i$ at the origin of each cycle.

The quasi-orthogonality of $\phi$ and $\mathcal{A}$ is supported by the
empirical observation that the temporal parameters $\{\bar{\Delta}^{+},
\bar{\Delta}^{-}, \bar{\phi}(H)\}$ have remained stable across three cycles
that exhibit dramatically different amplitude profiles ($130\times$, $20\times$,
$8\times$). The clock of the system does not slow down as the amplitude
compresses: the period has shortened by only $18$ days across three cycles even
as the multiplier contracted by a factor of sixteen. This decoupling is
precisely what makes the temporal dimension the more tractable and operationally
reliable forecasting variable.

\subsection{The Harmonic Model as a Formal Statement}

We now state the Unified Harmonic Model in a compact form that integrates
the empirical parameters derived across Sections~2--4.

\begin{definition}[Unified Harmonic Model]
The Bitcoin Harmonic Cycle is a two-dimensional dynamic system
$\mathbf{S}(t) = (\phi(t), \mathcal{A}(t))$ governed by the following
structural equations:

\medskip
\textit{Temporal law (period and internal phase symmetry):}
\begin{align}
  \bar{T}_{\text{cycle}} &= \bar{\Delta}^{+} + \bar{\Delta}^{-}
    = 1{,}060 + 370 = 1{,}430 \text{ days} \label{eq:period_law} \\[4pt]
  \bar{\phi}(H) &= \frac{\bar{\alpha}}{\bar{\alpha} + \bar{\beta}}
    = \frac{523}{1{,}059} \approx 0.494
    \quad \bigl(\varepsilon_\phi \leq 0.012\bigr) \label{eq:phase_law}
\end{align}

\textit{Amplitude law (power-law with ETF regime adjustment):}
\begin{align}
  \hat{\mathcal{M}}(d) &= 328 \cdot d^{\,5.1}
    \quad \text{(pre-ETF calibration)} \label{eq:amp_law_pre} \\[4pt]
  \mathcal{M}_i^{\text{adj}} &\in [2.5,\; 4.0]
    \quad \text{(regime-adjusted, Cycles } \geq 4\text{)} \label{eq:amp_law_etf}
\end{align}
\end{definition}

The model contains five structural parameters: two temporal means
($\bar{\Delta}^{+}$, $\bar{\Delta}^{-}$), one phase anchor
($\bar{\phi}(H)$), and two amplitude parameters ($\hat{\kappa} = 328$,
$\hat{\gamma} = 5.1$). All five are identified from historical data alone,
without recourse to price-level models, on-chain metrics, or macro-financial
variables. This parsimony is a deliberate methodological choice: the fewer the
free parameters, the higher the out-of-sample credibility of the projections
the model generates.

\subsection{Interaction Between Temporal and Amplitude Dimensions}

Although quasi-orthogonal in their dynamics, the temporal and amplitude
dimensions interact through two channels that are relevant for forward
projection.

\textbf{Channel I: Amplitude decay does not disturb the temporal clock.}
The evidence to date suggests that the progressive compression of the amplitude
envelope does not feed back into the period of the system. This is consistent
with the damped harmonic oscillator analogy: in equation~\eqref{eq:dho}, the
damping ratio $\zeta$ modifies the amplitude envelope $e^{-\zeta\omega_0 t}$
and the damped frequency $\omega_d = \omega_0\sqrt{1-\zeta^2}$, but for
$\zeta \ll 1$ the frequency shift is negligible. The empirical stability
$\mathcal{S}(\Delta^{+}) = 0.983$ obtained in Section~2 is consistent with
a regime of weak damping in the temporal dimension, even as the amplitude
damping appears considerably stronger. This decoupling is the single most
valuable operational property of the model.

\textbf{Channel II: ETF-driven drawdown compression creates a non-linear
feedback into amplitude.} As argued in Section~4, the reduction in prior
drawdown $d_i$ produced by structural ETF demand feeds into the amplitude
law through the steep exponent $\hat{\gamma} = 5.1$. A drawdown reduction
from $0.48$ to $0.35$---a change of 13 percentage points---translates into a
multiplier reduction from $8.1\times$ to approximately $2.5\times$: a
compression factor of more than three on the return side. This non-linear
amplification of amplitude consequences from relatively small changes in
drawdown depth is the primary risk to amplitude-dependent price projections
and the principal reason we present Cycle~4 amplitude estimates as a range
rather than a point forecast.

\subsection{The Calendar as the Primary Forecasting Instrument}

A corollary of the quasi-orthogonality result is of practical importance for
the evaluation of the model: temporal predictions can be confirmed or falsified
\textit{independently of price outcomes}. If the cycle minimum of 2026 arrives
within the projected window---regardless of whether the price reaches
\$35,000 or \$50,000---the temporal structure of the model survives. If it
does not, the model fails on its most robust dimension. This separation between
the testability of the temporal and amplitude predictions is what elevates the
framework beyond a price target exercise and positions it as a structural
hypothesis about the dynamical architecture of the market.


% =============================================================================
\section{Forward Projections as Falsifiable Hypotheses}
% =============================================================================

\subsection{Methodological Preamble}

The projections presented in this section are derived mechanically from the
parameters of the Unified Harmonic Model developed in Sections~2--5. They are
not price targets in the conventional analytical sense. They are
\textit{falsifiable hypotheses}: dated predictions whose confirmation or
refutation will determine the empirical viability of the model. Each projection
is accompanied by an explicit uncertainty interval that reflects the maximum
inter-cycle deviation observed in the corresponding historical parameter.

Price ranges are provided as supplementary estimates and must be understood as
secondary outputs of the model, dependent on the more uncertain amplitude
dimension and subject to the ETF regime-adjustment caveat of
Section~\ref{sec:etf_break}. The temporal projections are primary; the price
projections are indicative.

\subsection{Milestone I --- Cycle 3 Peak ($P_3$)}

The peak of the current expansion phase is projected at:

\begin{equation}
  \hat{P}_3 = 6\ \text{October}\ 2025
              \pm 45\ \text{days}
  \label{eq:p3_projection}
\end{equation}

This date is derived from two independent paths that converge on the same
interval:

\begin{enumerate}
  \item \textit{From the trough:} $T_3 = 21$ November 2022 $+ \bar{\Delta}^{+}
  = 1{,}060\ \text{days} \Rightarrow$ 16 October 2025.
  \item \textit{From the halving:} $H_3 = 20$ April 2024 $+ \bar{\beta} =
  536\ \text{days} \Rightarrow$ 8 October 2025.
\end{enumerate}

The two paths converge to within two days of each other, providing strong
internal consistency. The uncertainty window of $\pm 45$ days---wider than the
$\pm 12$ day parameter uncertainty of $\bar{\beta}$---is set conservatively to
account for the compounding of temporal uncertainty in the trough identification
and the possibility of a marginal new high being established slightly outside
the primary parameter range.

\medskip
\noindent\textit{Indicative price range:} \$135,000 -- \$160,000.

This range is derived by applying the Cycle 3 multiplier
$\mathcal{M}_3 \approx 8.1\times$ (under the hypothesis $P_3$ = \$108,000
equivalent at the assumed ATH level) and bounding it by the uncertainty in
the trough price identification. It is presented for contextual reference only
and does not constitute a price forecast.

\subsection{Milestone II --- Cycle 4 Trough ($T_4$)}

The minimum of the next contraction phase is projected at:

\begin{equation}
  \hat{T}_4 = \hat{P}_3 + \bar{\Delta}^{-}
            \approx 6\ \text{October}\ 2025 + 370\ \text{days}
            = 11\ \text{October}\ 2026
            \pm 45\ \text{days}
  \label{eq:t4_projection}
\end{equation}

The $\pm 45$ day window encompasses the full historical range of the
contraction phase ($363$ to $376$ days, a span of $13$ days) and an additional
buffer reflecting the dependency of this projection on the confirmation of
$\hat{P}_3$. The projection window therefore spans approximately from
late August 2026 to late November 2026.

\medskip
\noindent\textit{Indicative price range:} \$35,000 -- \$45,000.

This range assumes an ETF-dampened drawdown from $\hat{P}_3$ of approximately
$65$--$70\%$, consistent with the observed secular compression of drawdowns
($83\% \to 59\% \to 48\%$) and the additional stabilising effect of
sustained institutional demand during the contraction phase. A drawdown of
$65$--$70\%$ from a peak in the \$135,000--\$160,000 range produces a trough
in approximately the \$40,000--\$56,000 interval, with the lower portion of
the stated range reflecting potential peak undershooting.

\subsection{Milestone III --- Halving $H_4$ (Protocol-Determined)}

The fourth Bitcoin halving is an exogenous event determined by the protocol's
block generation schedule. Under the assumption of a sustained mean block time
of approximately 10 minutes, the halving at block height 1{,}050{,}000 is
expected at:

\begin{equation}
  H_4 \approx \text{April}\ 2028
  \label{eq:h4}
\end{equation}

Unlike all other milestones in this section, $H_4$ is not a model projection
but a protocol-derived datum. Its uncertainty is entirely a function of the
variance in block production rates and is, for practical purposes, contained
within a window of a few weeks around the stated month. This makes $H_4$ the
single fixed coordinate of the Cycle~4 temporal structure and the most reliable
anchor for downstream projections.

\subsection{Milestone IV --- Cycle 4 Peak ($P_4$)}

The peak of the next expansion phase admits two projection paths, both of which
converge on the second half of 2029:

\begin{enumerate}
  \item \textit{From the trough (full expansion):}
  \begin{equation}
    \hat{P}_4^{(a)} = \hat{T}_4 + \bar{\Delta}^{+}
                    \approx 11\ \text{October}\ 2026 + 1{,}060\ \text{days}
                    \approx 5\ \text{September}\ 2029
    \label{eq:p4_a}
  \end{equation}

  \item \textit{From the halving (post-halving segment):}
  \begin{equation}
    \hat{P}_4^{(b)} = H_4 + \bar{\beta}
                    \approx \text{April}\ 2028 + 536\ \text{days}
                    \approx \text{October}\ 2029
    \label{eq:p4_b}
  \end{equation}
\end{enumerate}

The two paths converge to within four to six weeks of each other, spanning the
interval from late August to late October 2029. We adopt as the central
projection:

\begin{equation}
  \hat{P}_4 = \text{Q3--Q4}\ 2029 \quad (\text{September--October}\ 2029)
              \pm 45\ \text{days}
  \label{eq:p4_projection}
\end{equation}

\medskip
\noindent\textit{Indicative price range:} \$120,000 -- \$180,000.

This range is derived from the regime-adjusted amplitude interval
$\mathcal{M}_4^{\text{adj}} \in [2.5, 4.0]$ applied to the projected trough
$\hat{T}_4^{\$} \in [\$35{,}000, \$45{,}000]$:

\begin{align}
  \text{Lower bound:} &\quad 2.5 \times \$35{,}000 = \$87{,}500 \\
  \text{Upper bound:} &\quad 4.0 \times \$45{,}000 = \$180{,}000
  \label{eq:p4_price}
\end{align}

The lower tail of the mechanical computation ($\approx\$87{,}500$) falls below
the stated indicative range. We adjust this lower bound upward to
\$120,000 on the grounds that a sub-ATH peak in Cycle~4 would represent a
structural break from the observed sequence of monotonically increasing nominal
peaks across all prior cycles---a break that would require independent
justification beyond the amplitude model. The \$120,000--\$180,000 range is
therefore the \textit{model-consistent} zone for $P_4$ under the assumptions
of the harmonic framework.

\subsection{Summary of Projections}

\begin{table}[H]
  \centering
  \caption{Forward Projection Summary --- Unified Harmonic Model}
  \label{tab:projections}
  \renewcommand{\arraystretch}{1.45}
  \begin{tabularx}{\textwidth}{l l l X}
    \toprule
    \textbf{Milestone} &
    \textbf{Projected Date} &
    \textbf{Price Range (ind.)} &
    \textbf{Derivation Basis} \\
    \midrule
    $\hat{P}_3$ --- Cycle 3 Peak
      & Oct 2025 $\pm$ 45d
      & \$135k -- \$160k
      & $T_3 + \bar{\Delta}^{+}$; $H_3 + \bar{\beta}$ \\[4pt]
    $\hat{T}_4$ --- Cycle 4 Trough
      & Oct 2026 $\pm$ 45d
      & \$35k -- \$45k
      & $\hat{P}_3 + \bar{\Delta}^{-}$; drawdown 65--70\% \\[4pt]
    $H_4$ --- Halving (protocol)
      & Apr 2028
      & ---
      & Block 1,050,000 (protocol-fixed) \\[4pt]
    $\hat{P}_4$ --- Cycle 4 Peak
      & Sep--Oct 2029 $\pm$ 45d
      & \$120k -- \$180k
      & $\hat{T}_4 + \bar{\Delta}^{+}$; $H_4 + \bar{\beta}$;
        $\mathcal{M}_4^{\text{adj}} \in [2.5, 4.0]$ \\
    \bottomrule
  \end{tabularx}
  \medskip
  \begin{minipage}{\textwidth}
    \small
    All temporal projections carry an explicit uncertainty of $\pm 45$ days.
    Price ranges are indicative and secondary to temporal projections.
    ``ind.'' = indicative. All projections constitute falsifiable hypotheses
    and must be revised upon material deviation from the projected temporal
    milestones.
  \end{minipage}
\end{table}


% =============================================================================
\section{Concluding Remarks and Model Limitations}
% =============================================================================

\subsection{The Calendar as Falsification Criterion}

This paper has advanced and formalised a single central thesis: the true
governing clock of Bitcoin's market cycles is not the halving event in
isolation, but a two-phase dynamic system in which the halving functions as a
temporal node---an internal anchor that divides the expansion phase into two
near-equal sub-intervals of approximately 523 and 536 days. The complete
expansion phase persists for approximately 1,060 days and is followed by a
contraction of approximately 370 days, yielding a mean total period of
1,430 days. The amplitude of each cycle, governed by a power-law relationship
with prior drawdown depth, declines monotonically and is subject to structural
acceleration under the ETF-driven institutional regime that emerged in 2024.

The principal intellectual contribution of this framework is not the price
targets it generates---those are secondary and heavily caveated---but rather
the methodological position it establishes: \textit{temporal structure is
the most robust and independently testable dimension of Bitcoin's market
dynamics.} Unlike price models, which can be retrospectively adjusted,
extended, or reinterpreted, the temporal projections of the Harmonic Model
are committed to specific calendar windows and are falsifiable without
ambiguity.

If the Cycle~4 trough does not materialise within the window of August--
November 2026, the temporal dimension of the model is falsified and requires
fundamental revision. If it does materialise within that window, the model
survives one additional out-of-sample test and its structural validity is
incrementally strengthened. The strength of a scientific model is not its
capacity to explain what has already occurred, but its willingness to be
specific about what has not yet occurred and to accept the consequences of
being wrong. The Harmonic Model is explicit on both counts.

\subsection{Model Limitations}

We conclude with a transparent assessment of the conditions under which the
model is expected to fail, partially or entirely.

\textbf{Limitation I: Small sample size.} The model is calibrated on three
complete cycles (with the third only provisionally closed). In classical
statistical terms, three observations are insufficient to establish the
significance of a structural parameter with high confidence. The temporal
stability indices derived in Section~2 are compelling but not formally
significant under conventional frequentist criteria. The model's credibility
is therefore primarily structural---based on the internal consistency of the
parameters and their convergence across independent measurement paths---
rather than statistical in the hypothesis-testing sense.

\textbf{Limitation II: Macro-financial black swan events.} The temporal
stability documented in this paper spans three cycles from 2015 to 2025. This
period encompassed considerable macroeconomic turbulence---negative interest
rates, a global pandemic, the fastest monetary tightening cycle in four
decades---without materially disrupting the temporal structure of Bitcoin's
market cycles. However, a sufficiently severe macro dislocation could alter
the clock. Scenarios that could produce structural disruption include: a
prolonged global recession accompanied by forced liquidation of risk assets
across all categories; a systemic failure of major custodial or ETF
infrastructure leading to forced outflows; or a sovereign-level regulatory
intervention that materially constrains the liquidity pool of the asset.
None of these scenarios is assigned a probability here; they are enumerated
as the class of exogenous shocks capable of falsifying the temporal dimension
of the model independently of its internal logic.

\textbf{Limitation III: Regulatory and structural regime change.} The
structural break introduced by spot ETF approval in 2024 has been incorporated
into the amplitude dimension of the model, but its full implications for the
temporal dimension remain uncertain. It is conceivable that sustained
institutional demand, by continuously absorbing selling pressure during
contraction phases, could either shorten or elongate the contraction duration
$\Delta^{-}$ relative to the historical average of 370 days. A systematic
shortening would be consistent with the hypothesis that institutional price
floors compress bear markets; a systematic elongation could arise if
institutional rebalancing generates prolonged gradual outflows rather than
abrupt capitulation. The direction and magnitude of this effect cannot be
determined from prior data and constitutes the most significant structural
uncertainty in the current cycle's contraction projection.

\textbf{Limitation IV: Protocol-level modifications.} The model treats the
four-year halving schedule as exogenous and immutable. Any modification to
Bitcoin's consensus rules that alters the block reward schedule---whilst
considered a negligible probability by the research and technical community at
the time of writing---would invalidate the halving's role as a temporal
node and require a complete reconstruction of the phase anchor parameter
$\bar{\phi}(H)$.

\subsection{Final Remark}

The Harmonic Time Model does not claim to predict the price of Bitcoin. It
claims something structurally prior and more durable: that the \textit{timing}
of Bitcoin's market cycles exhibits a degree of internal regularity that has
persisted across radically different market regimes, capitalisation levels, and
structural environments---and that this regularity is sufficiently precise to
generate dated, falsifiable predictions about future cycle milestones.

The model's value will be determined not by the elegance of its mathematical
formulation, but by the accuracy of the projection entered in
Table~\ref{tab:projections}. If the trough of Cycle~4 arrives in October 2026,
and the peak of Cycle~4 arrives in the third or fourth quarter of 2029, the
Harmonic Model will have earned the right to be taken seriously as a
structural description of Bitcoin's dynamic architecture. Until then, it
remains what all empirical models should be: a precise and refutable
hypothesis about the behaviour of a complex system.

% =============================================================================
%  END OF SECTIONS 5, 6 & 7
%  (The \end{document} tag belongs in the base template file.)
% =============================================================================


% =============================================================================
%  BIBLIOGRAPHY (Failsafe Empty Setup)
% =============================================================================
\begin{thebibliography}{99}
\bibitem{nakamoto2008}
Nakamoto, S. (2008). \textit{Bitcoin: A Peer-to-Peer Electronic Cash System}. Self-published.
\bibitem{romero2026}
Romero Quant. (2026). \textit{The Invariance of Bitcoin's Internal Expansion Phases}. Quantitative Research Series, Vol. 1.
\end{thebibliography}

\end{document}
% =============================================================================
%  END OF BASE TEMPLATE
% =============================================================================